% 原书物理页=印刷页12--19；第2章开头 + 2.1 全节. 2.1.10 末尾的注与习题位于 p20 上部（2.2 标题之前），仍属 2.1，一并译出.
\chapter{Lie 代数：概述}\label{chap:ssl-2}
\section{模与表示}\label{sec:ssl-2-1}
设 $K$ 是一个交换域. 把结合 $K$-代数上模的定义照搬过来，就得到下面的概念.
\begin{definition}\label{def:ssl-2-1-1}
$K$ 上Lie 代数 $\mathfrak{g}$ 上的一个\newterm{ssl-module}{模}{module}（或称\newterm{ssl-g-module}{$\mathfrak{g}$-模}{g-module}）是指一个 $K$-向量空间 $V$，它带有一个 $K$-双线性映射
\[\begin{gathered}
\mathfrak{g} \times V \longrightarrow V,\\
(X, v) \longmapsto X \cdot v
\end{gathered}\]
使得对一切 $X, Y \in \mathfrak{g}$ 与 $v \in V$ 都有 $[X, Y] \cdot v = X \cdot (Y \cdot v) - Y \cdot (X \cdot v)$.
\end{definition}
换言之，若以 $r(X) \in \operatorname{End}_K(V)$ 记自同态 $v \mapsto X \cdot v$，则二元组 $(V, r)$ 就是前文所给意义下 $\mathfrak{g}$ 的一个表示. 因此应当把表示看成结合代数上模的类似物. 由此我们常常简化记号，把 $r(X)(v)$ 写作 $X \cdot v$（甚至 $Xv$），并且不加区分地使用``$\mathfrak{g}$-模''或``$\mathfrak{g}$ 的表示''这两种说法.
\begin{example*}
伴随表示不是别的，正是让 $\mathfrak{g}$ 通过括号 $X \cdot Y := [X, Y]$ 左作用于自身而得到的 $\mathfrak{g}$-模 $\mathfrak{g}$.
\end{example*}

\begin{numpara}[$\mathfrak{g}$ 的表示范畴]\label{np:ssl-2-1-2}
与结合代数的情形一样，$\mathfrak{g}$-模之间的态射是与 $\mathfrak{g}$ 的作用交换的线性映射. 用表示的语言说，$(V, r)$ 与 $(V', r')$ 之间的态射是满足 $\forall X \in \mathfrak{g},\ r'(X) \circ f = f \circ r(X)$ 的 $K$-线性映射 $f : V \longrightarrow V'$；用简化记号即 $\forall X \in \mathfrak{g},\ \forall v \in V,\ f(Xv) = Xf(v)$. 我们记 $\operatorname{Hom}_{\mathfrak{g}}(V, V')$ 为由这些态射构成的 $K$-向量空间.
这样得到一个范畴 $\operatorname{Rep}_K(\mathfrak{g})$（或简记 $\operatorname{Rep}(\mathfrak{g})$），其对象是表示，态射如上. 它其实是一个 Abel 范畴，并且核、余核与直和都同到 $K$-向量空间范畴的遗忘函子交换. 换言之，若 $f : V \longrightarrow V'$ 是表示之间的态射，那么它的核（按向量空间的意义）是 $V$ 的子表示（或者说 $V$ 的 $\mathfrak{g}$-子模），它的像是 $V'$ 的子表示，而它的余核具有唯一的表示结构，使得从 $V'$ 出发的投影与作用相容（``商表示''）. 类似地，两个表示的直和定义为其底层空间的直和，配上 $\mathfrak{g}$ 的显然作用.
\begin{example*}
$\mathfrak{g}$ 的理想不是别的，正是 $\mathfrak{g}$ 的 $\mathfrak{g}$-子模，亦即 $\mathfrak{g}$ 的伴随表示的一个子表示.
\end{example*}
我们还把 $\operatorname{Rep}(\mathfrak{g})$ 中由 $K$ 上有限维表示构成的满子范畴记作 $\operatorname{Rep}^{\mathrm{df}}(\mathfrak{g})$.
\end{numpara}

\begin{numpara}[Lie 群的表示与``微分''函子]\label{np:ssl-2-1-3}
$\mathfrak{g}$-模的概念源自结合代数的语境. 之所以同时也使用``表示''这一术语，是因为它们与 Lie 群之间有紧密的联系. 在本段中假设 $K = \mathbb{R}$ 或 $\mathbb{C}$.
若 $(V, \rho)$ 是 $G$ 的一个表示，就可以通过微分给它配备 $\operatorname{Lie}(G)$ 的一个表示 $(V, d\rho)$. 此外，若 $f : V \longrightarrow V'$ 是 $G$ 的表示之间的态射，那么它也是相应的 $\operatorname{Lie}(G)$ 表示之间的态射，因为在点 $g = e$ 处对等式 $f \circ \rho(g) = \rho'(g) \circ f$ 求导，便得 $\forall X \in \operatorname{Lie}(G),\ f \circ d\rho(X) = d\rho'(X) \circ f$. 这样就得到一个正合函子
\[\begin{gathered}
\operatorname{Rep}^{\mathrm{df}}_K(G) \longrightarrow \operatorname{Rep}^{\mathrm{df}}_K(\operatorname{Lie}(G)),\\
(V, \rho) \longmapsto (V, d\rho)
\end{gathered}\]
它是研究 $G$ 的表示时最重要的工具.
\begin{proposition*}
上述``微分''函子：
\begin{enumerate}[label=\roman*)]
\item 当 $G$ 连通时是全忠实的；
\item 当 $G$ 单连通时是本质满的（从而是范畴等价）.
\end{enumerate}
\end{proposition*}
参见 [GALM2], 1.12.2.
\begin{remark*}
这里重要的是限于有限维表示. $\mathfrak{sl}_2(\mathbb{C})$ 存在无法积分到 $\mathrm{SL}_2(\mathbb{C})$ 的无限维表示.
\end{remark*}
下面的一些定义与构造，同关于群表示的定义与构造是平行的（其中大多数在结合代数上的模那边没有对应物）. 因此，把两种语境并行处理是有意思的.
\end{numpara}

\begin{numpara}[平凡表示、单位表示与不变向量]\label{np:ssl-2-1-4}
设 $V$ 是群 $G$ 或Lie 代数 $\mathfrak{g}$ 的某个表示的向量空间. 记
\begin{itemize}[label=\textendash]
\item （群）$V^{G} := \{v \in V,\ \forall g \in G,\ gv = v\}$，即由 $G$ 下不变的向量构成的子空间；
\item （代数）$V^{\mathfrak{g}} := \{v \in V,\ \forall X \in \mathfrak{g},\ Xv = 0\}$，即由被 $\mathfrak{g}$ 零化的向量构成的子空间.
\end{itemize}
\begin{remark*}
若 $(V, \rho)$ 是连通 Lie 群 $G$ 的表示，容易证明 $V^{\operatorname{Lie}(G)} = V^{G}$.
\end{remark*}
群 $G$（相应地，Lie 代数 $\mathfrak{g}$）的表示 $V$ 称为\newterm{ssl-trivial-representation}{平凡表示}{trivial representation}，如果 $V = V^{G}$（相应地，$V = V^{\mathfrak{g}}$）. $G$（相应地，$\mathfrak{g}$）的\newterm{ssl-unit-representation}{单位表示}{unit representation}是指带有平凡表示的向量空间 $K$；在无歧义时简单地记作 $K$.
\end{numpara}

\begin{numpara}[张量积、逆步表示与内部 Hom]\label{np:ssl-2-1-5}
设 $V, V'$ 是某个群的表示 $\rho, \rho'$（或某个Lie 代数 $\mathfrak{g}$ 的表示 $r, r'$）的底层 $K$-向量空间.
\begin{enumerate}[label=\roman*)]
\item 我们在 $V \otimes_K V'$ 上赋予\newterm{ssl-tensor-product}{张量积}{tensor product}表示，其定义为：
\begin{itemize}[label=\textendash]
\item （群）$g \mapsto \rho(g) \otimes \rho'(g)$，亦即 $g(v \otimes v') := gv \otimes gv'$，$\forall v \in V,\ v' \in V'$；
\item （代数）$X \mapsto r(X) \otimes \mathrm{id}_{V'} + \mathrm{id}_V \otimes r'(X)$，亦即 $X(v \otimes v') := Xv \otimes v' + v \otimes Xv'$，$\forall v \in V,\ v' \in V'$.
\end{itemize}
\begin{exercise*}
验证第二个定义确实给出 $\mathfrak{g}$ 的一个表示，再验证当 $G$ 是 Lie 群时，$\rho \otimes \rho'$ 的微分正是 $d\rho \otimes \mathrm{id} + \mathrm{id} \otimes d\rho'$.
\end{exercise*}
\begin{exercise*}
若 $V' = K$ 是单位表示，验证 $V \otimes_K K \simeq V$（这也解释了``单位表示''这一术语）.
\end{exercise*}
\begin{remark*}
对任意的结合代数 $A$，并不存在 $A$-模张量积的概念.
\end{remark*}
\item 我们在 $\operatorname{Hom}_K(V, V')$ 上赋予\newterm{ssl-hom-interne}{内部 $\mathrm{Hom}$}{internal Hom}表示，其定义为：
\begin{itemize}[label=\textendash]
\item （群）$(g, f) \mapsto \rho'(g) \circ f \circ \rho(g^{-1})$，亦即 $(gf)(v) := g\,f(g^{-1}v)$，$\forall v \in V$；
\item （代数）$(X, f) \mapsto r'(X) \circ f - f \circ r(X)$，亦即 $(Xf)(v) := Xf(v) - f(Xv)$，$\forall v \in V$.
\end{itemize}
我们把这个表示记作 $\mathcal{H}om_K(V, V')$，称为内部 $\mathrm{Hom}$（其中 $H$ 为花体）.
\begin{exercise*}
验证这些公式确实定义了表示，并且对 Lie 群 $G$ 而言，群的内部 $\mathrm{Hom}$ 表示的微分正是 $\operatorname{Lie}(G)$ 的内部 $\mathrm{Hom}$ 表示.
\end{exercise*}
\begin{remark*}
有 $\operatorname{Hom}_G(V, V') = \mathcal{H}om_K(V, V')^{G}$ 与 $\operatorname{Hom}_{\mathfrak{g}}(V, V') = \mathcal{H}om_K(V, V')^{\mathfrak{g}}$.
\end{remark*}
\item ii) 的一个重要特例是 $V' = K$ 为单位表示的情形. 这时简单地记 $V^{\vee} := \mathcal{H}om_K(V, K)$，并称之为 $V$ 的\newterm{ssl-contragredient}{逆步表示}{contragredient representation}（也称 $V$ 的对偶表示）.
\end{enumerate}
张量积与内部 $\mathrm{Hom}$ 按下述方式相联系：
\begin{lemma*}
若 $V$ 与 $V'$ 都是有限维，则态射 $V^{\vee} \otimes_K V' \longrightarrow \mathcal{H}om_K(V, V')$（它把 $v^{*} \otimes v'$ 映为态射 $(v, v') \mapsto \langle v, v^{*}\rangle\, v'$\footnote{译注：原文此处作 $(v, v')$；按文义该态射应为 $v \mapsto \langle v, v^{*}\rangle\, v'$.}）是表示的同构.
\end{lemma*}
\begin{proof}
留作习题.
\end{proof}
\end{numpara}

\begin{numpara}[表示的约化]\label{np:ssl-2-1-6}
（群或Lie 代数的）表示 $V$ 称为：
\begin{itemize}[label=\textendash]
\item \newterm{ssl-irreducible}{不可约}{irreducible}的，如果它仅有的子表示是 $\{0\}$ 与 $V$；
\item \newterm{ssl-completely-reducible}{完全可约}{completely reducible}的，如果它是不可约子表示的直和；
\item \newterm{ssl-indecomposable}{不可分解}{indecomposable}的，如果它不是两个子表示的直和.
\end{itemize}
\begin{example*}
（$\mathfrak{g} = K$ 的情形）设 $\mathfrak{g}$ 是一维的，取 $X \in \mathfrak{g}$. $\mathfrak{g}$ 的表示 $(V, r)$ 于是由自同态 $r(X) \in \operatorname{End}_K(V)$ 确定. 当 $K$ 代数闭时，可以看出：
\begin{itemize}[label=\textendash]
\item $V$ 不可约当且仅当 $\dim_K(V) = 1$（存在特征直线.）；
\item $V$ 完全可约当且仅当 $r(X)$ 可对角化；
\item $V$ 不可分解当且仅当在 Jordan 分解 $r(X) = r(X)_s + r(X)_n$ 中，$r(X)_n$ 的幂零指数等于 $V$ 的维数（此时 $r(X)_s$ 是一个位似）. 换句话说，$r(X)$ 只有一个 Jordan 块.
\end{itemize}
然而，当 $K$ 不代数闭时，$\mathfrak{g}$ 的不可约表示的维数可以大于 $1$. 例如当 $K = \mathbb{R}$ 时，由
\[r(X) = \begin{pmatrix} 0 & 1 \\ -1 & 0 \end{pmatrix}\]
给出的二维表示.
\end{example*}
\begin{remark*}
若 $(V, \rho)$ 是连通 Lie 群 $G$ 的有限维表示，则 $(V, \rho)$ 不可约、不可分解或完全可约，当且仅当 $(V, d\rho)$ 相应地如此.
\end{remark*}
\begin{exercise*}
证明 $V$ 不可约、不可分解或完全可约当且仅当 $V^{\vee}$ 相应地如此.
\end{exercise*}
对维数作递降归纳，容易看出每个有限维表示都是不可分解子表示的直和. 因此，要理解范畴 $\operatorname{Rep}^{\mathrm{df}}(G)$ 或 $\operatorname{Rep}^{\mathrm{df}}(\mathfrak{g})$，先验地需要对不可分解表示进行分类. 在上面的例子中，它们由维数与 $X$ 的唯一特征值分类. 但一般而言，这个问题可能变得过于困难、甚至毫无希望，例如对Abel（即括号为零的）Lie 代数 $\mathfrak{g} = K^3$（它的一个表示相当于给定 $V$ 上三个两两交换的自同态）.
如果知道 $G$ 或 $\mathfrak{g}$ 的表示都完全可约，事情就大大简化. 这时，应用下面的 Schur 引理可知只需对不可约表示分类，而后者往往更容易.
\end{numpara}

\begin{lemma}[Schur]\label{lemma:ssl-2-1-7}
设 $f : V \longrightarrow V'$ 是不可约表示之间的态射. 若 $f$ 非零，则 $\operatorname{Ker}(f) \neq V$，于是由 $V$ 的不可约性得 $\operatorname{Ker}(f) = 0$；同理 $\operatorname{Im}(f) \neq 0$，由 $V'$ 的不可约性得 $\operatorname{Im}(f) = V'$. 因此 $f$ 要么为零，要么可逆.
特别地，若 $V$ 是有限维不可约表示，则 $\operatorname{End}_{\mathfrak{g}}(V)$ 是一个有限维的 $K$-可除代数. 当 $K$ 代数闭时，这蕴涵 $\operatorname{End}_{\mathfrak{g}}(V) = K$（Schur 引理）.
\end{lemma}
\textbf{应用.}\quad 完全可约的表示 $V$ 可以唯一地（不计次序）写成 $V = V_1^{\oplus e_1} \oplus \cdots \oplus V_r^{\oplus e_r}$ 的形式，其中 $V_1, \cdots, V_r$ 是两两不同构的不可约表示. 事实上，重数 $e_i$ 由公式
\[e_i = \dim_K(\operatorname{Hom}_{\mathfrak{g}}(V_i, V)) / \dim_K(\operatorname{End}_{\mathfrak{g}}(V_i))\]
给出.
\begin{remark*}
对于不完全可约的表示，也存在一个关于不可分解表示直和分解的唯一性结果，即所谓 Krull--Schmidt--Azumaya 定理. 其证明更为精细，我们将用不到它. 感兴趣的读者可参看 Curtis 与 Reiner 的表示论著作之一.
\end{remark*}

\begin{numpara}\label{np:ssl-2-1-8}
我们以后会看到，特征为零的域上的``半单''Lie 代数满足这条完全可约性. 在此之前，先让我们熟悉这一概念.
\begin{lemma*}
对表示 $V$，下列性质等价：
\begin{enumerate}[label=\roman*)]
\item $V$ 的任何子表示 $W$ 都有稳定补（从而是一个子表示）$W'$ 使 $W' \oplus W = V$；
\item $V$ 完全可约，
\item $V$ 是它的不可约子表示之和.
\end{enumerate}
\end{lemma*}
\begin{proof}
i)$\Rightarrow$ii). 对 $V$ 的维数作归纳. 若 $V$ 不可约（特别地，若 $\dim(V) = 1$），则无事可做. 否则，取 $V$ 的一个维数 $> 0$ 最小的子表示 $W$. 它是不可约的，且由假设有稳定补 $V'$. 我们证明 $V'$ 仍满足 i). 设 $W' \subset V'$ 是一个子表示，取 $W \oplus W'$ 在 $V$ 中的稳定补 $W''$. 则 $W''$ 沿 $W$ 到 $V'$ 上的投影是 $W'$ 在 $V'$ 中的稳定补. 于是 $V'$ 满足 i)，由归纳假设，$V'$ 是不可约表示的直和. 故 $V = W \oplus V'$ 也是.
ii)$\Rightarrow$iii) 是同义反复.
iii)$\Rightarrow$i). 证明 $W \subset V$ 有稳定补，对 $W$ 的余维数作归纳. 若余维数为零，则无事可做. 否则，存在（由假设）$V$ 的不可约表示 $U$ 使得 $U \not\subset W$，从而 $W \cap U = \{0\}$. 对 $W \oplus U$ 应用归纳假设，得到 $W \oplus U$ 的一个稳定补 $W'_1$，于是只须令 $W' := U \oplus W'_1$.
\end{proof}
下面给出完全可约性成立的情形的一个例子.
\begin{example*}
设 $K = \mathbb{C}$. 设 $V$ 是带有一个 Hermite 内积的表示，内积按下面的意义不变：
\begin{itemize}[label=\textendash]
\item （群）$\forall v, v' \in V,\ \forall g \in G,\ \langle gv, gv'\rangle = \langle v, v'\rangle$；
\item （代数）$\forall v, v' \in V,\ \forall X \in \mathfrak{g},\ \langle Xv, v'\rangle + \langle v, Xv'\rangle = 0$.
\end{itemize}
则 $V$ 完全可约. 事实上，若 $W \subset V$ 是一个稳定子空间，则它的正交 $W^{\perp} = \{v \in V,\ \forall w \in W,\ \langle v, w\rangle = 0\}$ 与 $W$ 成直和 $W \oplus W^{\perp}$（因为 $\langle\,\cdot\,,\,\cdot\,\rangle$ 是正定的）并且是稳定的.
\end{example*}
\begin{exercise*}
设 $G$ 是有限群，$V$ 是 $G$ 的表示，$\langle\,\cdot\,,\,\cdot\,\rangle$ 是 $V$ 上的一个 Hermite 内积. 证明公式 $(v, w) := \frac{1}{|G|}\sum_{g \in G} \langle gv, gw\rangle$ 定义一个 $G$ 不变的 Hermite 内积. 由此推出有限群的每个复表示都完全可约.
\end{exercise*}
\end{numpara}

\begin{numpara}[积的表示]\label{np:ssl-2-1-9}
设 $G, G'$ 是两个群，或 $\mathfrak{g}, \mathfrak{g}'$ 是两个Lie 代数. 若 $W$ 是 $G$（相应地，$\mathfrak{g}$）的复表示，$W'$ 是 $G'$（相应地，$\mathfrak{g}'$）的表示，则在 $W \otimes_K W'$ 上赋予 $G \times G'$（相应地，$\mathfrak{g} \oplus \mathfrak{g}'$）的表示：$(g, g')(v \otimes v') := (gv, gv')$\footnote{译注：原文此处作 $(gv, gv')$；按文义应为 $gv \otimes g'v'$（对照下文代数情形的公式）.}，相应地 $(X, X')(v \otimes v') := Xv \otimes v' + v \otimes X'v'$.
\begin{proposition*}
设 $K$ 代数闭. 在（群或代数）两种情形的每一种中，都有：
\begin{enumerate}[label=\roman*)]
\item 若 $W$ 与 $W'$ 不可约，则 $W \otimes_K W'$ 也不可约（\emph{对积而言！}）；
\item $G \times G'$（相应地，$\mathfrak{g} \oplus \mathfrak{g}'$）的每个不可约表示 $V$ 都具有上述形状 $W \otimes_K W'$；
\item 若 $W$ 与 $W'$ 完全可约，则 $W \otimes_K W'$ 亦然.
\end{enumerate}
\end{proposition*}
\end{numpara}
\begin{proof}
我们处理Lie 代数的情形；群的情形完全类似. 先做一个一般的构造. 若 $W \in \operatorname{Rep}(\mathfrak{g})$，$V \in \operatorname{Rep}(\mathfrak{g} \oplus \mathfrak{g}')$，我们在 $\operatorname{Hom}_{\mathfrak{g}}(W, V)$ 上赋予 $\mathfrak{g}'$ 的表示，定义为 $(X'f)(w) := X' \cdot f(w)$（这是``内部 $\mathrm{Hom}$''构造的一个变体）. 于是有赋值映射
\[\begin{gathered}
ev_{W,V} : W \otimes_K \operatorname{Hom}_{\mathfrak{g}}(W, V) \longrightarrow V,\\
w \otimes f \longmapsto f(w).
\end{gathered}\]
容易验证它是 $\mathfrak{g} \oplus \mathfrak{g}'$ 的表示之间的态射（习题：请验证！）.
现在设 $V$ 与 $W$ 都不可约，且 $W$ 同构于 $V|_{\mathfrak{g}}$ 的一个不可约子表示. 此时 $ev_{W,V}$ 非零，从而是满射. 设 $n' := \dim_K(\operatorname{Hom}_{\mathfrak{g}}(W, V))$. 于是有满射 $W^{\oplus n'} \twoheadrightarrow V$，它表明 $V$ 作为 $\mathfrak{g}$-模完全可约，并且同构于某个 $W^{\oplus e}$，其中 $e \leqslant n'$. 而此时，由于 $K$ 代数闭，Schur 引理告诉我们 $n' = e$，于是最终证明了 $ev_{W,V}$ 是同构. 由此还得出 $W' := \operatorname{Hom}_{\mathfrak{g}}(W, V)$ 是不可约 $\mathfrak{g}'$-模，并且作为 $\mathfrak{g}'$-模有 $V \simeq W'^{\oplus n}$，其中 $n = \dim_K(W)$.
于是我们证明了 ii). 把它应用于 $W \otimes_K W'$ 的一个非零不可约子表示，又可得到 i).
剩下证明 iii). 但若 $W = \bigoplus_i W_i$，$W' = \bigoplus_j W'_j$，其中诸 $W_i$ 与诸 $W'_j$ 不可约，则 $W \otimes W' = \bigoplus_{i,j} W_i \otimes W'_j$ 就是不可约表示之和的一个分解.
\end{proof}

\begin{numpara}[特征为零时 $\mathfrak{g} = \mathfrak{sl}_2$ 的不可约表示]\label{np:ssl-2-1-10}
作为例子——但更主要的是由于这里用到的技巧与中间结果对我们日后研究``半单''Lie 代数至关重要——我们在``$K$ 是\newterm{ssl-alg-closed-char0-field}{特征为零的代数闭域}{algebraically closed field of characteristic zero}''的假设下分类 $\mathfrak{sl}_2$ 的不可约表示.
回忆 $\mathfrak{sl}_2$ 是 $3$ 维的. 我们将使用下面的基 $(E, H, F)$：
\[E = \begin{pmatrix} 0 & 1 \\ 0 & 0 \end{pmatrix},\quad H = \begin{pmatrix} 1 & 0 \\ 0 & -1 \end{pmatrix},\quad F = \begin{pmatrix} 0 & 0 \\ 1 & 0 \end{pmatrix}.\]
$\mathfrak{g} := \mathfrak{sl}_2$ 的括号由下列公式确定：
\[[H, E] = 2E,\qquad [H, F] = -2F,\qquad [E, F] = H.\]
若 $(V, r)$ 是 $\mathfrak{g} = \mathfrak{sl}_2$ 的表示且 $\lambda \in K$，记
\[V_{\lambda} := \operatorname{Ker}(r(H) - \lambda\, \mathrm{id}_V).\]
若 $V_{\lambda} \neq 0$，它就是 $r(H)$ 的属于特征值 $\lambda$ 的特征子空间. 这时称 $\lambda$ 是 $H$ 在 $V$ 中的一个\newterm{ssl-weight}{权}{weight}.
\begin{lemma*}
设 $(V, r)$ 是 $\mathfrak{g}$ 的表示，$\lambda \in K$ 是 $H$ 在 $V$ 中的一个权. 则 $r(E)V_{\lambda} \subset V_{\lambda+2}$ 且 $r(F)V_{\lambda} \subset V_{\lambda-2}$.
\end{lemma*}
\begin{proof}
设 $v \in V_{\lambda}$. 有 $HEv = EHv + [H, E]v = \lambda Ev + 2Ev = (\lambda + 2)Ev$. 同理 $HFv = (\lambda - 2)Fv$.
\end{proof}
\begin{proposition*}
对每个整数 $m \in \mathbb{N}$，在同构意义下存在唯一的 $m + 1$ 维不可约表示 $(V_m, r_m)$. 此外，$H$ 在 $V_m$ 中的权集为 $\{-m, -m + 2, \cdots, m - 2, m\}$.
\end{proposition*}
\begin{proof}
i) 先证唯一性. 设 $V$ 是 $\mathfrak{g}$ 的有限维不可约表示. 由于 $K$ 代数闭，$V$ 的权集非空. 又因这个集合有限而 $K$ 特征为零，故存在权 $\lambda_0$ 使 $\lambda_0 + 2$ 不是权. 取非零的 $v_0 \in V_{\lambda_0}$. 由前一引理，此时 $EV_{\lambda_0} = 0$.
置 $v_k := \tfrac{1}{k!}F^k v_0$，并令 $W = \operatorname{Vect}_K\{v_k,\ k \in \mathbb{N}\}$ 为诸 $v_k$ 张成的空间. 按构造，这个空间在 $F$ 下稳定；又因诸 $v_k$ 是 $H$ 的特征向量（前一引理），它在 $H$ 下也稳定. 再证它在 $E$ 下稳定. 已见 $Ev_0 = 0$. 对 $k > 1$，有
\[Ev_k = \frac{1}{k} EFv_{k-1} = \frac{1}{k} \bigl( FEv_{k-1} + [E, F] v_{k-1} \bigr) = \frac{1}{k} \bigl( FEv_{k-1} + (\lambda_0 - 2k + 2) v_{k-1} \bigr).\]
由此可用归纳法证明 $Ev_k = (\lambda_0 - k + 1) v_{k-1}$.
由此得出 $W$ 是 $V$ 的非零子表示，再由不可约性得 $W = V$. 诸非零的 $v_k$ 线性无关，因为它们对应于 $H$ 的两两不同的特征值；于是令 $m = \dim_K(V) - 1$，便得 $V$ 的一组基 $v_0, \cdots, v_m$.
为确定 $\lambda_0$，注意由 $H = [E, F]$ 必有 $\operatorname{tr}(r(H)) = 0$. 而 $\operatorname{tr}(r(H)) = \lambda_0 + (\lambda_0 - 2) + \cdots + (\lambda_0 - 2m) = (m + 1)\lambda_0 - m(m + 1)$，解得 $\lambda_0 = m$；最终，诸权恰如所述.
ii) 现在须验证存在性. 为此只须验证 $M_{m+1}(K)$ 中如下的矩阵：
\[H_m := \sum_{i=1}^{m+1} (m - 2i + 2) E_{i,i},\qquad E_m := \sum_{i=1}^{m} (m - i + 1) E_{i,i+1},\qquad \text{及}\ F_m := \sum_{i=1}^{m} i E_{i+1,i}\]
满足关系 $[H_m, E_m] = 2E_m$、$[H_m, F_m] = -2F_m$ 与 $[E_m, F_m] = H_m$. 此后只须置 $r_m(H) := H_m$、$r_m(E) := E_m$ 与 $r_m(F) := F_m$.
\end{proof}
\begin{remark*}
在低维时我们认出一些熟悉的表示：
\begin{itemize}[label=\textendash]
\item $(V_0, r_0)$ 是 $\mathfrak{g}$ 的单位表示；
\item $(V_1, r_1)$ 是由包含 $\mathfrak{g} \subset \mathfrak{gl}_2(K)$ 给出的``标准''表示；
\item $(V_2, r_2)$ 是伴随表示 $\mathfrak{g} \longrightarrow \mathfrak{gl}(\mathfrak{g}_K) \simeq \mathfrak{gl}_3(K)$.
\end{itemize}
一般地，可以证明 $(V_m, r_m) \simeq \operatorname{Sym}^m(V_1, r_1)$，即 $(V_1, r_1)$ 的 $m$ 次对称幂，其底层空间为 $\operatorname{Sym}^m(V_1) = (V_1 \otimes V_1 \otimes \cdots \otimes V_1)^{\mathfrak{S}_m}$，而 $\mathfrak{g}$ 的作用是张量积表示的商.
\end{remark*}
\begin{remark*}
表示 $(V_m, r_m)$ 称为\newterm{ssl-highest-weight}{最高权}{highest weight}为 $m$ 的不可约表示. 注意，证明中真正用到的基 $E, H, F$ 的性质只有公式 $[H, E] = 2E$、$[H, F] = 2F$\footnote{译注：原文此处作 $[H, F] = 2F$；按本节开头的括号公式应为 $[H, F] = -2F$，疑漏负号.}与 $[E, F] = H$. 任何满足这些性质的基 $(E', H', F')$ 在 $\mathrm{GL}_2(K)$ 的作用下共轭于 $(E, H, F)$. 事实上，把上述证明中 ii) 的公式应用于标准表示，就给出 $K^2$ 的一组基 $v_0, v_1$，使得 $(E', H', F')$ 在这组基下的矩阵恰是 $E, H, F$.
\end{remark*}
\begin{exercise*}
\begin{enumerate}[label=\roman*)]
\item 证明由 $r(H) = x\,\frac{\partial}{\partial x} - y\,\frac{\partial}{\partial y}$、$r(E) = x\,\frac{\partial}{\partial y}$ 与 $r(F) = y\,\frac{\partial}{\partial x}$ 定义的 $K$-线性映射 $\mathfrak{sl}_2(K) \xrightarrow{\ r\ } \operatorname{End}_K(K[x, y])$ 是 $\mathfrak{sl}_2(K)$ 的一个表示.
\item 设 $\mathcal{P}_m := \{f(x, y) \in K[x, y],\ \text{$m$ 次齐次}\}$. 证明 $\mathcal{P}_m$ 在 $\mathfrak{sl}_2(K)$ 下稳定，且如此得到的表示 $(\mathcal{P}_m, r)$ 同构于表示 $(V_m, r_m)$，即最高权为 $m$ 的不可约表示.
\end{enumerate}
\end{exercise*}
\end{numpara}
